\section{Data separation and inference units}
E2 crosses two geometry schemes, source rasters 48/64 and five paired seeds
(20 models), with 800 training and 100 validation scenes per condition. Training
follows Section~\ref{sec:task}; the source rasters share RGB and the final
checkpoint is fixed without validation selection. Each independent-geometry
test bank has 400 scenes. Training uses one assignment version per scene;
evaluation retains both but samples one version/card for the actual decision.

E3 preselects all five independent/source 64 E2 checkpoints. It uses 400 development
scenes and, per domain, 1,600 calibration and 2,000 test scenes. The normal
generator and independent path sampler in Section~\ref{sec:task} operate before
rasterization; no pixel, risk or feasibility screen changes their sampling.

Oracle discovery reuses the consumed E3 IID/appearance test banks. The later
prospective bank has 2,000 new IID scenes (geometry seed 2612091301), with a recorded
source/checkpoint freeze before generation. Exact anonymous patch and complete-
scene identities are disjoint from all 10,200 E2/E3 records. IDs, appearance names,
seeds and path/patch ordering are excluded from the identity comparison; arbitrary
geometric symmetries are not quotient out. After that confirmation, the bank is
consumed. The PLIC/shift analysis in Section~\ref{sec:comparison} is explicitly
appended to it, without new samples or renewed confirmatory status.

Learned-track ratios pool five models over the same 2,000 scenes; crossed
bootstrap draws preserve that dependence. Target-oracle intervals resample
scenes only. Five training seeds still limit the approximate crossed inference.

\section{Spatial reference and reconstruction details}
The direct reference rasterizes each continuous property and swept-union footprint
at 512. It averages their product over footprint area, then applies the public
card law. It is a fixed discrete reference for synthetic mean exposure. The
auxiliary property check rasterizes properties at 1024, averages to 512 and retains
F512. An earlier full 256/direct 512 comparison changes both property and footprint
rasters and is kept separate from the common-footprint attribution.

\paragraph{Signed cost attribution.}
Exposure replacement follows Eq.~\ref{eq:telescope}. For every fixed action, cost
differences are separately evaluated as
\[
\widehat c-\refcost=(\widehat c-\targetcost)
 +(\targetcost-\poolcost)+(\poolcost-\refcost).
\]
The nonlinear formula is applied before differencing. Channel exposures,
signed/absolute costs and label-pattern transitions are preserved. The variance
envelope uses privileged reference occupancies; its zero-flip implications are
implementation checks.

\paragraph{PLIC area inversion.}
The normal is the negative weighted centered gradient of the coarse fraction,
with transverse weights $(1,2,1)$ and edge-replicated padding. Pairing differences
before summation eliminates roundoff-generated transverse components for an
exactly constant transverse field. Normalize $(n_x,n_y)$ by its L1 norm. For a
nonzero normal let $ a=|n_x|$, $ b=|n_y|$, $ a+b=1$, $ s=\min(a,b)$ and $ l=\max(a,b)$.
The area CDF of $ aU+bV$ for independent unit-square coordinates can be evaluated
symmetrically. At $0\le t\le 1/2$ its lower half is
\[
H(t)=\begin{cases}t^2/(2ab),&t<s,\\(t-s/2)/l,&t\ge s.\end{cases}
\]
Axis-aligned limits are handled analytically; the upper half uses $1-H(1-t)$.
For fraction $ p $, invert the lower tail $ q=\min(p,1-p)$ using
$ t=\sqrt{2abq}$ when $ q\le s/(2l)$, otherwise $ t=lq+s/2$; reflect for $ p>1/2$ and
translate for negative normal components. This yields a half-plane whose
unit-square area is the original $ p $.

For each fine pixel square, the same CDF gives its exact fractional area under
the reconstructed half-plane. Integrating those fractions against binary F512
preserves each coarse cell's original mass. Cells where a footprint is entirely
zero or one need no fine calculation: the integral is fixed by conserved mass.
All fractions, including small learned probabilities, are retained. A mixed
cell whose gradient L1 norm is at most $10^{-12}$ stays uniform. This is an
explicit fallback, not a inferred true boundary.

On the appended bank, the maximum analytic inverse-area discrepancy is
$2.78\times 10^{-16}$ and the maximum sum-of-fine-areas discrepancy is
$3.33\times 10^{-16}$, below the fixed $10^{-10}$ tolerance. There are 34,837
gradient-fallback cells across the input tracks and 5,047,140 reconstructed query
cells. These numerical checks establish conservation, not physical correctness
of the inferred boundary or a safety guarantee.

\section{Shift fitting and unequal risk points}
\label{app:shift}
The fixed cost-offset bank is
\[
\{- .02,- .015,- .01,- .005,- .002,- .001,0,.001,.002,.005,.01,.015,.02\}.
\]
One offset is fitted for perfect target input and one is shared across the five
model inputs. Fitting uses only consumed IID discovery data: maximize safe
acceptance subject to unsafe-accepted count no larger than uniform. Ties prefer
smallest absolute offset and then smaller offset. The rule is frozen at .001
for target input and 0 for learned input. It is empirical fitting, not a risk
certificate or guarantee that the same unsafe-count constraint will hold later.

The entire development bank is retained. At target offset .005, development has
1,578 safe and 17 unsafe accepts; at .01 it has 1,598 safe and 33 unsafe accepts.
Fixed PLIC has 1,600 safe and 16 unsafe accepts on development. The selected .001
shift has 1,550 safe and one unsafe accept. Larger shifts therefore can recover
more opportunities, and the selected point must not be interpreted as the
capacity of the whole shift family. No alternative offset is chosen using the
appended outcomes. The different risk levels preclude claiming a matched-risk
deployment advantage from the reported selected points.

\begin{table}[h]
\centering\small
\caption{Absolute appended outcomes. Learned rows pool five models sharing
2,000 scenes. The direct oracle has 1,592 safe opportunities per model.}
\begin{tabular}{llrrrr}
\toprule
Input & Interpretation & Accepted & Safe accepted & Unsafe accepted & Safe opportunities lost\\
\midrule
Target & Uniform &1540&1538&2&54\\
 & Shift &1549&1544&5&48\\
 & PLIC &1600&1582&18&10\\
Learned & Uniform &7783&7703&80&257\\
 & Shift &7783&7703&80&257\\
 & PLIC &7909&7786&123&174\\
\bottomrule
\end{tabular}
\end{table}

\section{Event-level cancellation evidence}
Each recovery, new safe loss, new unsafe accept and resolved unsafe accept keeps
its scene/model/card/action identity. At the original uniform-selected action,
write target-measurement error as $ m=c_{\rm target}-\refcost $ and model-relative
error as $ l=\widehat c-c_{\rm target}$. After changing interpretation, retain both
$ m'$ and $ l'$, with
$\widehat c'-\widehat c=(m'-m)+(l'-l)$. This fixed-action equality is checked
before interpreting reselected outcomes.

For the 43 newly unsafe learned PLIC decisions, 25 satisfy all of: unchanged
action, $ ml<0$ before modification, $|m'|<|m|$, and
$|m'+l'|>|m+l|$. The other 18 are retained as counterevidence to an exclusive
cancellation explanation. The first matching event in scene/model order is
scene 264, model index 1, exchanged version, card D, action 0. In it, both $|m|$
and $|l|$ decrease yet total underestimation grows, as detailed in the main text.
This is a transparent conditional illustration; it was not chosen to estimate
the population frequency of such cases.

\section{Prospective direction check and its addendum}
\label{app:direction}
The original confirmation uses one-sided budgets $.05/3$. H1 counts scenes where
direct 512 and the finer-property check both retain at least one safe candidate,
but pooled 512 fails to produce a safe accepted action. Its 61/2,000 events give
exact lower .02286937. H2 initially uses a crossed bootstrap for the fraction of
fixed-selected pooling flips that are conservative. Since all 322 flips are
conservative, its bootstrap lower degenerates to 1 and cannot establish population
certainty. H3 compares small/larger margin flip rates under joint boundary
overlap; its estimate is .373080 and approximate lower .260586.

The direction addendum was recorded at 20:57:54 UTC after job submission but
before any fresh outcome was inspected and before the original hypothesis file
was generated at 21:00:24 UTC. It is disclosed as a strengthening after the
original freeze, not silently included in that freeze. It requires an additional
exact scene-level check at the same H2 error budget: conditional on a scene
having any flip across the fixed five-model bank, every flip must be conservative.
All 68 changed scenes satisfy it; the exact lower is .941566. This conditional
fixed-bank estimand differs from the model--scene weighted proportion. Both
original outputs and the added check are preserved.

\section{Earlier controls and closed external attempt}
\label{app:earlier}
E1 uses five historical models and a consumed 300-scene calibration bank. With
restored output, positive-residual quantiles are .02797--.03287 for all 24
pair/version/card/candidate entries, .02201--.02767 for the current four
candidates, and .00501--.00790 for the fixed selected action. These are three
different protection objects. The smaller marginal selected-action bound does
not itself control unsafe frequency after acceptance.

E3 separately uses all five preselected E2 independent/source 64 models. Twelve
threshold rules per source/target calibration domain use exact one-sided
Clopper--Pearson upper bounds with family budget
$\delta/(5\cdot 2\cdot 12)=.05/120$. The desired conditional unsafe limit is
$\epsilon=.05$, distinct from danger threshold $\tau=.02$ and residual miscoverage
$\alpha=.05$. The original fixed .02 restored rules already certify, with source
upper bounds 1.66--2.54\%; selecting a higher-coverage certified threshold increases
both opportunity use and test risk. Source certificates do not become unknown-
distribution guarantees on appearance shift, and no old certificate applies to
PLIC or a newly fitted shift.

The RELLIS attempt evaluated a fixed 24-frame training-sequence pilot. All-four
corridor unknown area at most 5\% held in 0/24 cases; an optimistic diagnostic
discarding depth/occlusion checks reached 3/24. These results concern the chosen
single-scan surface-fit interface, not the dataset's general suitability. A
bounded revision inspected 240 raw scans around those same 24 development frames
and ran a motion-compensation prototype. Independent physical alignment and
visibility error bounds remained unresolved. It stopped before selecting,
acquiring or evaluating the planned new 48-frame set. This is a prerequisite
stop, not 0/48 and not a mechanism-transfer result; neither external training nor
an external final test was run.

\section{Computational verification}
The original oracle round has six distinct tests and independent replay of
52,000 discovery decisions and 26,000 prospective decisions, with 768 and 384
blockwise candidate-cost checks. The additional PLIC tests cover mass against
independent polygon clipping, uniform-lift identity, zero-gradient fallback,
complement/transposition behavior, axis orientation, and constant-footprint
cells. One initial test detected transverse roundoff at approximately
$8.3\times 10^{-16}$ for an exactly constant transverse field; changing the
arithmetic to pair differences before summing fixed it. The original test,
failed attempt, and unchanged scientific tolerance are retained; no development
scene was read before the repaired tests passed.

Independent appended verification replays 72,000 decisions and 65,536 fine-square
areas with polygon clipping. The maximum polygon difference is
$3.33\times10^{-16}$; the constant-lift identity differs by at most
$3.71\times10^{-14}$. It also verifies all 25 specified cancellation-pattern
events among the 43 new unsafe learned accepts. All computation runs on the cluster; frozen inputs and outputs are preserved.
